\recipeTitle{Pasta del Patatrack Hackathon}

\recipeSubtitle{Una ricetta segnaposto per iniziare il ricettario collaborativo.}

\recipeAuthor{Felice Pantaleo and contributors}

\recipeStory{Questa e' una ricetta di esempio pensata per mostrare il formato standard del progetto. In futuro potra' essere sostituita da una vera ricetta nata durante un Patatrack hackathon.}

% Add a real image in recipes/patatrack-hackathon-pasta/images/ and uncomment the line below.
% \recipePhoto{recipes/patatrack-hackathon-pasta/images/patatrack-hackathon-pasta.jpg}

\recipeInfo
  {10 minuti}
  {12 minuti}
  {4 persone}
  {Facile}
  {Primo piatto}

\begin{ingredients}
  \ingredient{Pasta}{400 g}
  \ingredient{Olio extravergine di oliva}{3 cucchiai}
  \ingredient{Aglio}{1 spicchio}
  \ingredient{Parmigiano grattugiato}{40 g}
  \ingredient{Sale}{q.b.}
\end{ingredients}

\begin{tools}
  \tool{Pentola grande}
  \tool{Padella}
  \tool{Coltello}
  \tool{Tagliere}
\end{tools}

\begin{procedure}
  \step{Porta a bollore abbondante acqua in una pentola grande, sala l'acqua e cuoci 400 g di pasta seguendo il tempo indicato sulla confezione.}
  \step{Nel frattempo scalda 3 cucchiai di olio extravergine di oliva in una padella con 1 spicchio d'aglio schiacciato.}
  \step{Scola 400 g di pasta tenendo da parte un mestolo di acqua di cottura.}
  \step{Salta 400 g di pasta nella padella con l'olio, aggiungendo poca acqua di cottura se serve per legare.}
  \step{Spegni il fuoco, elimina 1 spicchio d'aglio e aggiungi 40 g di parmigiano grattugiato. Mescola e servi subito.}
\end{procedure}

\begin{notes}
  \note{Questa ricetta e' volutamente semplice: serve come template pratico per nuove ricette.}
\end{notes}
